\documentclass[journal]{IEEEtran}
\usepackage{cite}
\usepackage{amsmath,amssymb,amsfonts}
\usepackage{algorithmic}
\usepackage{graphicx}
\usepackage{textcomp}
\usepackage{xcolor}
\usepackage{url}
\usepackage{booktabs}
\usepackage{multirow}

\begin{document}

\title{Supplementary Material for "Unobtrusive Web-Based Sensing Framework for Longitudinal Affective Monitoring"}

\maketitle

\section{FaceMesh Landmark Indices}
The MediaPipe FaceMesh model returns 468 landmarks. For calculating the metrics used in this study, we utilized the specific indices listed in Table \ref{tab:landmarks}.
The "Rigid Alignment" indices were selected based on their low variance during facial expression (i.e., they move with the head, not the skin).

\begin{table}[h]
\caption{Selected Facial Landmark Indices}
\label{tab:landmarks}
\centering
\begin{tabular}{ll}
\bottomrule
\textbf{Region of Interest} & \textbf{Mesh Indices (MediaPipe)} \\ 
\toprule
\textbf{Nose Tip (tracking)} & 1 \\
\textbf{Left Tragus (rigid)} & 234, 127, 162 \\
\textbf{Right Tragus (rigid)} & 454, 356, 389 \\
\textbf{Left Eye (EAR)} & 33, 160, 158, 133, 153, 144 \\
\textbf{Right Eye (EAR)} & 362, 385, 387, 263, 373, 380 \\
\textbf{Forehead (rPPG ROI)} & 10, 338, 297, 332, 284, 251, 389, 356 \\
\bottomrule
\end{tabular}
\end{table}

\section{Signal Processing Specifications}
The raw physiological signals extracted from the vision pipeline are inherently noisy. 
We applied the following digital signal processing (DSP) chain using the WebAudio API's `AudioContext` for real-time filtering, or `scipy.signal` for post-processing.

\subsection{rPPG Bandpass Filter}
To isolate the cardiac frequency band from motion artifacts and DC offsets:
\begin{itemize}
    \item \textbf{Type}: 4th-order Butterworth Filter (Infinite Impulse Response).
    \item \textbf{High-Pass Cutoff}: 0.7 Hz (42 BPM). Removes slow drifts due to shadow movement.
    \item \textbf{Low-Pass Cutoff}: 3.0 Hz (180 BPM). Removes 60Hz power line noise and sensor thermal noise.
    \item \textbf{Zero-Phase Implementation}: Applied bidirectionally (`filtfilt`) to prevent phase distortion which would alter HRV timing.
\end{itemize}

\subsection{Motion Smoothing}
The raw head coordinate $P_t$ was smoothed using an Exponential Moving Average (EMA) to reduce quantization noise from the webcam sensor:
\begin{equation}
    S_t = \alpha \cdot P_t + (1 - \alpha) \cdot S_{t-1}
\end{equation}
Where $\alpha = 0.2$, equivalent to a time constant of $\tau \approx 80ms$ at 60 FPS. 
This preserves the sharp onset of a "Freeze" response while smoothing out jitter.

\section{System Latency Analysis}
A critical requirement for psychophysics is timing precision. 
We profiled the "End-to-End" latency of the Camerala system using a high-speed camera (240 FPS) observing the screen-to-photodiode transition.

\begin{enumerate}
    \item \textbf{Input Latency}: $8.4ms \pm 2.1$ (USB Keyboard via `keydown` event).
    \item \textbf{Display Latency}: $16.7ms$ (Fixed by 60Hz V-Sync).
    \item \textbf{Processing Latency}: $12.3ms \pm 4.5$ (MediaPipe Inference Time on Intel Iris Plus Graphics).
\end{enumerate}
Total System Latency is approximately $37.4ms$. 
Since the Reaction Time (RT) analysis searches for differences $>50ms$ between conditions, this latency floor is acceptable.

\end{document}
